\subsection{Preesquemas geométricamente unirramosos.}
\label{IV.2.s6.15-fr}

\begin{env}[6.15.1]
\label{IV.2.s6.15.1-fr}
Diremos que un preesquema $X$ es \emph{unirramoso} (resp.
\emph{geométricamente unirramoso}) en un punto $x$, o que el
punto $x$ es \emph{unirramoso} (resp. \emph{geométricamente unirramoso}) en $X$,
si el anillo local $\mathcal O_{X,x}$ es \emph{unirramoso} (resp.
\emph{geométricamente unirramoso}) (\textbf 0, 23.2.1). Diremos que $X$ es
\emph{unirramoso} (resp. \emph{geométricamente unirramoso}) si lo es en cada
uno de sus puntos. Se sigue de esta definición que, para que $X$
sea unirramoso (resp. geométricamente unirramoso) en un punto,
es necesario y suficiente que $X_{\mathrm{red}}$ lo sea en ese
punto.
\end{env}

\begin{env}[6.15.2]
\EGAReaderTarget{ega4.c04.statement.6-15-2}{IV.2.s6.15.2-fr}
Hay que tener presente que, con las definiciones de (6.15.1),
\emph{un anillo local} $A$ puede ser unirramoso (resp. geométricamente
unirramoso) \emph{sin que} $\operatorname{Spec}(A)$ \emph{lo sea}. En otras
palabras, puede haber ideales primos $\mathfrak p$ de $A$ tales
que $A_\mathfrak p$ no sea unirramoso; esto equivale también a
decir que, en un preesquema, la noción de punto geométricamente
unirramoso \emph{no es estable por generización}. Lo veremos en el
ejemplo siguiente.
\end{env}
